\title{\textbf{Applied Vacuum Engineering} \\ \large \textit{Volume I: Foundations and Universal Operators}}
\author{Grant Lindblom}
\date{}

\maketitle

\vfill
\noindent\textbf{Applied Vacuum Engineering: Volume I} \\
This volume establishes the mechanical and mathematical axioms of the $\mathcal{M}_A$ LC Condensate framework, derives the universal operators ($Z$, $S$, $\Gamma$), and defines the four-regime map that governs all subsequent volumes.

\begin{abstract}
    The Standard Model requires more than 26 independent empirical parameters. Applied Vacuum Engineering (AVE) models the vacuum as a structured medium and reduces those inputs to three calibration boundaries: the spatial cutoff ($\ell_{node}$), the dielectric yield bound ($\alpha$), and the macroscopic gravitational coupling ($G$).
    
    \textbf{Volume I: Foundations and Universal Operators} explicitly defines the continuum architecture of the vacuum, derives the universal operators $Z = \sqrt{\mu/\epsilon}$, $S = \sqrt{1 - r^2}$, and $\Gamma = (Z_2 - Z_1)/(Z_2 + Z_1)$, and establishes the four-regime map that every subsequent volume builds upon.
    \begin{equation}
        Z_0 = \sqrt{\frac{L}{C}} \quad , \quad c = \frac{1}{\sqrt{\epsilon_0 \mu_0}}
    \end{equation}
    On this basis the framework develops quantum mechanics (the Born rule, entanglement, and the generalized uncertainty principle) as emergent, deterministic signal dynamics on the substrate rather than as primitive postulates.
\end{abstract}

